\chapter{Hito 6: Construcción del Primer Proceso ETL con DuckDB}

\section{Objetivo}

El objetivo de este hito es construir el primer flujo completo de carga de datos utilizando DuckDB.

Por primera vez conectaremos:

\begin{verbatim}
CSV
|
V
Python
|
V
DuckDB
\end{verbatim}

Al finalizar este hito el estudiante será capaz de:

\begin{itemize}
\item Leer archivos CSV mediante Pandas.
\item Conectarse a DuckDB desde Python.
\item Cargar información a la capa staging.
\item Automatizar el proceso de carga.
\item Validar que los datos fueron cargados correctamente.
\end{itemize}

\section{Resultado esperado}

Al finalizar este hito los datos contenidos en:

\begin{verbatim}
data/raw/tipo_cambio.csv

data/raw/balances.csv
\end{verbatim}

deberán encontrarse cargados en:

\begin{verbatim}
staging.tipo_cambio_raw

staging.balance_bancario_raw
\end{verbatim}

dentro de DuckDB.

\section{Arquitectura del proceso}

El flujo construido será:

\begin{verbatim}
CSV
|
V
Pandas
|
V
DuckDB
\end{verbatim}

\section{Paso 1: Verificar archivos fuente}

Confirmar la existencia de:

\begin{verbatim}
data/raw/tipo_cambio.csv

data/raw/balances.csv
\end{verbatim}

\section{Paso 2: Verificar conexión}

Verificar:

\begin{verbatim}
etl/db.py
\end{verbatim}

Contenido:

\begin{lstlisting}[language=Python]
import duckdb

DB_PATH = (
"data/warehouse/"
"macrofinanzas.duckdb"
)

def get_connection():

```
return duckdb.connect(
    DB_PATH
)
```

\end{lstlisting}

\section{Paso 3: Crear script de carga}

Crear:

\begin{verbatim}
etl/load_staging.py
\end{verbatim}

\section{Paso 4: Importar librerías}

Agregar:

\begin{lstlisting}[language=Python]
import pandas as pd

from db import get_connection
\end{lstlisting}

\section{Paso 5: Leer tipo de cambio}

Agregar:

\begin{lstlisting}[language=Python]
fx = pd.read_csv(
"data/raw/tipo_cambio.csv"
)

print(
f"Filas FX: {len(fx)}"
)
\end{lstlisting}

\section{Paso 6: Leer balances}

Agregar:

\begin{lstlisting}[language=Python]
balances = pd.read_csv(
"data/raw/balances.csv"
)

print(
f"Filas balances: {len(balances)}"
)
\end{lstlisting}

\section{Paso 7: Agregar metadatos}

Agregar:

\begin{lstlisting}[language=Python]
fx["archivo_origen"] = (
"tipo_cambio.csv"
)

balances["archivo_origen"] = (
"balances.csv"
)
\end{lstlisting}

Agregar timestamp:

\begin{lstlisting}[language=Python]
from datetime import datetime

fx["cargado_en"] = (
datetime.now()
)

balances["cargado_en"] = (
datetime.now()
)
\end{lstlisting}

\section{Paso 8: Conectarse a DuckDB}

Agregar:

\begin{lstlisting}[language=Python]
conn = get_connection()
\end{lstlisting}

\section{Paso 9: Limpiar staging}

Agregar:

\begin{lstlisting}[language=Python]
conn.execute(
"""
DELETE FROM
staging.tipo_cambio_raw;
"""
)

conn.execute(
"""
DELETE FROM
staging.balance_bancario_raw;
"""
)
\end{lstlisting}

Esta estrategia garantiza que la carga sea reproducible.

\section{Paso 10: Registrar DataFrames}

DuckDB permite consultar directamente DataFrames.

Agregar:

\begin{lstlisting}[language=Python]
conn.register(
"fx_df",
fx
)

conn.register(
"balances_df",
balances
)
\end{lstlisting}

\section{Paso 11: Cargar tipo de cambio}

Agregar:

\begin{lstlisting}[language=Python]
conn.execute(
"""
INSERT INTO
staging.tipo_cambio_raw

SELECT *
FROM fx_df;
"""
)
\end{lstlisting}

\section{Paso 12: Cargar balances}

Agregar:

\begin{lstlisting}[language=Python]
conn.execute(
"""
INSERT INTO
staging.balance_bancario_raw

SELECT *
FROM balances_df;
"""
)
\end{lstlisting}

\section{Paso 13: Verificar conteos}

Agregar:

\begin{lstlisting}[language=Python]
result_fx = conn.execute(
"""
SELECT COUNT(*)
FROM staging.tipo_cambio_raw
"""
).fetchone()[0]

result_balances = conn.execute(
"""
SELECT COUNT(*)
FROM staging.balance_bancario_raw
"""
).fetchone()[0]

print(
f"Filas staging FX: {result_fx}"
)

print(
f"Filas staging balances: {result_balances}"
)
\end{lstlisting}

\section{Paso 14: Cerrar conexión}

Agregar:

\begin{lstlisting}[language=Python]
conn.close()
\end{lstlisting}

\section{Versión completa del script}

\begin{lstlisting}[language=Python]
import pandas as pd

from datetime import datetime

from db import get_connection

fx = pd.read_csv(
"data/raw/tipo_cambio.csv"
)

balances = pd.read_csv(
"data/raw/balances.csv"
)

fx["archivo_origen"] = (
"tipo_cambio.csv"
)

balances["archivo_origen"] = (
"balances.csv"
)

fx["cargado_en"] = (
datetime.now()
)

balances["cargado_en"] = (
datetime.now()
)

conn = get_connection()

conn.execute(
"""
DELETE FROM
staging.tipo_cambio_raw;
"""
)

conn.execute(
"""
DELETE FROM
staging.balance_bancario_raw;
"""
)

conn.register(
"fx_df",
fx
)

conn.register(
"balances_df",
balances
)

conn.execute(
"""
INSERT INTO
staging.tipo_cambio_raw

SELECT *
FROM fx_df;
"""
)

conn.execute(
"""
INSERT INTO
staging.balance_bancario_raw

SELECT *
FROM balances_df;
"""
)

result_fx = conn.execute(
"""
SELECT COUNT(*)
FROM staging.tipo_cambio_raw
"""
).fetchone()[0]

result_balances = conn.execute(
"""
SELECT COUNT(*)
FROM staging.balance_bancario_raw
"""
).fetchone()[0]

print(
f"Filas staging FX: {result_fx}"
)

print(
f"Filas staging balances: {result_balances}"
)

conn.close()

print("Carga completada.")
\end{lstlisting}

\section{Paso 15: Ejecutar ETL}

Desde terminal:

\begin{lstlisting}[language=bash]
python etl/load_staging.py
\end{lstlisting}

Salida esperada:

\begin{verbatim}
Filas FX: 100
Filas balances: 50

Filas staging FX: 100
Filas staging balances: 50

Carga completada.
\end{verbatim}

\section{Paso 16: Crear consulta de monitoreo}

Crear:

\begin{verbatim}
sql/monitor_staging.sql
\end{verbatim}

Contenido:

\begin{lstlisting}[language=SQL]
SELECT
'tipo_cambio_raw' AS tabla,
COUNT(*) AS total_filas
FROM staging.tipo_cambio_raw

UNION ALL

SELECT
'balance_bancario_raw' AS tabla,
COUNT(*) AS total_filas
FROM staging.balance_bancario_raw;
\end{lstlisting}

\section{Paso 17: Ejecutar consulta}

Desde Python:

\begin{lstlisting}[language=Python]
result = conn.execute(
open(
"sql/monitor_staging.sql"
).read()
).fetchdf()

print(result)
\end{lstlisting}

Resultado esperado:

\begin{verbatim}
tabla  total_filas
0      tipo_cambio_raw          100
1  balance_bancario_raw           50
\end{verbatim}

\section{Paso 18: Crear orquestador}

Crear:

\begin{verbatim}
etl/run_pipeline.py
\end{verbatim}

Contenido inicial:

\begin{lstlisting}[language=Python]
import subprocess

subprocess.run(
[
"python",
"etl/load_staging.py"
],
check=True
)

print(
"Pipeline ejecutado correctamente."
)
\end{lstlisting}

\section{Buenas prácticas}

Durante este hito seguiremos las siguientes reglas:

\begin{enumerate}
\item Todas las cargas deben ser reproducibles.
\item No modificar tablas staging manualmente.
\item Mantener metadatos de carga.
\item Validar conteos después de cada carga.
\item Cerrar conexiones al finalizar.
\end{enumerate}

\section{Checklist de validación}

\begin{itemize}
\item[$\Box$] CSV leídos correctamente.
\item[$\Box$] Conexión DuckDB funcional.
\item[$\Box$] Tablas staging limpiadas.
\item[$\Box$] Datos cargados.
\item[$\Box$] Conteos validados.
\item[$\Box$] Script monitor_staging.sql ejecutado.
\item[$\Box$] Pipeline reproducible.
\end{itemize}

\section{Entregable}

Al finalizar este hito el estudiante deberá disponer de:

\begin{itemize}
\item Primer proceso ETL funcional.
\item Datos cargados en staging.
\item Consulta de monitoreo.
\item Orquestador inicial.
\item Flujo completo CSV → Python → DuckDB.
\end{itemize}

\section{Próximo hito}

En el Hito 7 construiremos la capa Core.

Crearemos:

\begin{itemize}
\item Dimensión Tiempo.
\item Dimensión Moneda.
\item Dimensión Entidad Financiera.
\item Claves subrogadas.
\item Primera transformación desde staging hacia el modelo dimensional.
\end{itemize}
